\scanpage{266}
\unsure{a)|(i)} $X$ est projectif sur $Y$

(ii) $\widehat X$ est projectif sur $\widehat Y$, i.e. $X\otimes_A\operatorname{Spec}(\widehat A)$ projectif sur $\operatorname{Spec}\widehat A$

(iii) Pour tout entier $n$, $X_n$ est projectif sur $Y_n$.

Évident $\text{a)}\Rightarrow\text{(ii)}\Rightarrow\text{(iii)}$. Prouvons $\text{(iii)}\Rightarrow\text{a)}$. Soit $n_0$ tel que l'image de $\operatorname{Pic}(X_{n_0})\unsure{-|\to}\operatorname{Pic}(X_0)$ soit universelle.

Soit $L_{n_0}$ Module inv. sur $X_{n_0}$ ample $/Y_{n_0}$. \deleted{\unsure{Il|Je} provient d'un $\unsure{L}$} \inserted{Donc $\unsure{L_0|L}$ en vertu de 2.1. (i)} provient d'un module inv. $\mathcal L$ sur $\widehat X$, donc en vertu de \deleted{\unsure{(i)}} 2.1. (ii) il provient aussi d'un module \inserted{inv.} $L$ sur $X$. Ce dernier est ample en vertu de EGA III. Qfd.

\margin{$L_0=L_{n_0}|X_0$
\figure{\begin{tikzpicture}[x=1cm,y=1cm]\draw (0,1)..controls(1,1.5)and(2,1.2)..(2.2,.5)--(2.15,-.1)--(.8,-.1)..controls(.35,-.2)and(.7,-.55)..(.1,-.5);\end{tikzpicture}}}

Notons maintenant que 2.1. \deleted{\unsure{(i)}} se déduit aisément de 1.7. par le \unsure{Gaga|Gag} de Serre-Grauert-Remmert-Grothendieck

\mspage{265}
